% Importado desde: 2026-03-28_Simetrias_Exactas_y_Grupo_Efectivo_del_Paso_Dirac.tex
\chapter{Derivacion Pedagogica: --- simetrias exactas del paso discreto --- y grupo efectivo emergente en Dirac \(3+1\)}
\label{ch:simetrias-grupo-paso-dirac}

\section{Objetivo}
En el capitulo anterior vimos que la construccion discreta tipo Dirac en \(3+1\) es consistente en el sector infrarrojo, pero globalmente vive sobre un objeto mas rico: toro de Brillouin, cuasienergia periodica, fibra interna de cuatro componentes y un paso unitario local ordenado.

Ahora toca fijar con cuidado una pregunta muy importante:
\begin{displayquote}
\textbf{cuales son exactamente las simetrias del paso discreto completo, y como se relacionan con el grupo continuo relativista que aparece solo en el limite infrarrojo?}
\end{displayquote}

Este capitulo quiere dejar clara la jerarquia correcta:
\begin{enumerate}[label=\arabic*.]
    \item un nivel \textbf{microscopico exacto}, con simetrias discretas del paso unitario;
    \item un nivel \textbf{efectivo infrarrojo}, donde emerge el Hamiltoniano de Dirac;
    \item y un nivel \textbf{espinorial relativista}, donde aparecen \(SO^+(3,1)\), Poincare y \(SL(2,\mathbb{C})\).
\end{enumerate}

\section{El objeto exacto}
Tomamos el bloque Weyl elemental
\begin{equation}
U_W(\vk)=
\ee^{-\,\ii a k_z \sigma_z}
\ee^{-\,\ii a k_y \sigma_y}
\ee^{-\,\ii a k_x \sigma_x},
\label{c25:eq:UW}
\end{equation}
y armamos la caminata Dirac de cuatro componentes como
\begin{equation}
U_D(\vk)=M\,U_0(\vk),
\qquad
M=\ee^{-\,\ii \Delta t\,m\,\tau_x\otimes \bbI},
\label{c25:eq:UD}
\end{equation}
con
\begin{equation}
U_0(\vk)=
\begin{pmatrix}
U_W(\vk) & 0\\
0 & U_{-W}(\vk)
\end{pmatrix},
\qquad
U_{-W}(\vk)=
\ee^{+\,\ii a k_x \sigma_x}
\ee^{+\,\ii a k_y \sigma_y}
\ee^{+\,\ii a k_z \sigma_z}.
\label{c25:eq:U0}
\end{equation}

En el infrarrojo,
\begin{equation}
U_D(\vk)\simeq \bbI-\ii\Delta t\,H_D,
\end{equation}
con
\begin{equation}
H_D=
v\,\tau_z\otimes(\sigma_x p_x+\sigma_y p_y+\sigma_z p_z)
+m\,\tau_x\otimes\bbI.
\label{c25:eq:HD}
\end{equation}

\section{Que significa una simetria exacta del paso}
Para una caminata discreta, la nocion correcta de simetria exacta no es \enquote{el Hamiltoniano parece isotropo}, sino la siguiente:

\begin{displayquote}
\textbf{Diremos que una transformacion espacial discreta \(g\) es simetria exacta del paso si existe una unitaria interna \(S_g\) tal que}
\[
S_g\,U_D(\vk)\,S_g^{-1}=U_D(g\vk)
\]
\textbf{para todo \(\vk\) de la zona de Brillouin.}
\end{displayquote}

Esta definicion separa bien dos cosas:
\begin{enumerate}[label=\arabic*.]
    \item la accion sobre el espacio de momentos;
    \item la accion espinorial sobre la fibra interna de cuatro componentes.
\end{enumerate}

Si la igualdad vale solo al expandir para \(\vk\approx 0\), entonces no es una simetria exacta, sino una \textbf{simetria efectiva}.

\section{Primer nivel: simetrias exactas garantizadas}
\subsection{Traslaciones espaciales discretas}
Como el mismo paso local actua en cada celda, la red admite invariancia exacta bajo traslaciones de \(\mathbb{Z}^3\). En espacio recíproco, eso se refleja en que \(\vk\) es una buena etiqueta de Bloch.

Luego la primera pieza del grupo microscopico es
\begin{equation}
\Lambda_{\mathrm{esp}}=\mathbb{Z}^3.
\end{equation}

\subsection{Traslacion temporal discreta}
El tiempo no es continuo en este nivel. La evolucion exacta viene dada por iteraciones enteras del paso:
\begin{equation}
\Psi_{n+1}=U_D\,\Psi_n.
\end{equation}
Por tanto, tambien hay una estructura exacta de traslaciones temporales discretas:
\begin{equation}
\Lambda_t=\mathbb{Z}.
\end{equation}

Esta observacion parece simple, pero es central: el grupo de tiempo exacto no es \(\mathbb{R}\), sino \(\mathbb{Z}\).

\section{Segundo nivel: simetrias puntuales exactas del paso minimo}
\subsection{La idea general}
Si el paso estuviera completamente simetrizado entre \(x\), \(y\) y \(z\), podriamos esperar un grupo puntual microscopico grande. Pero nuestro modelo minimo usa un producto \textbf{ordenado}:
\[
U_W=U_zU_yU_x.
\]

Ese orden deja memoria microscopica. Por eso no debemos asumir que todo el grupo cubico actua exactamente.

\subsection{Rotacion exacta de orden dos alrededor del eje z}
Consideremos
\begin{equation}
R_z(\pi):(k_x,k_y,k_z)\mapsto(-k_x,-k_y,k_z).
\end{equation}
La accion espinorial correspondiente sobre el bloque Weyl es
\begin{equation}
V_z=\ii\sigma_z.
\end{equation}

Usando
\begin{equation}
V_z\sigma_xV_z^{-1}=-\sigma_x,
\qquad
V_z\sigma_yV_z^{-1}=-\sigma_y,
\qquad
V_z\sigma_zV_z^{-1}=+\sigma_z,
\label{c25:eq:Vzaction}
\end{equation}
obtenemos
\begin{align}
V_z\,U_W(\vk)\,V_z^{-1}
&=
\ee^{-\,\ii a k_z \sigma_z}
\ee^{+\,\ii a k_y \sigma_y}
\ee^{+\,\ii a k_x \sigma_x}\notag\\
&=
U_W\!\bigl(R_z(\pi)\vk\bigr).
\label{c25:eq:C2zW}
\end{align}

El punto importante es que el orden \(z,y,x\) se conserva. Solo cambian signos.

\subsection{Extension al bloque Dirac}
En cuatro componentes tomamos
\begin{equation}
S_z=\bbI_{\tau}\otimes V_z.
\end{equation}
Entonces
\begin{equation}
S_z\,U_0(\vk)\,S_z^{-1}=U_0\!\bigl(R_z(\pi)\vk\bigr).
\end{equation}
Ademas, como el termino de masa es proporcional a \(\tau_x\otimes\bbI\), conmuta con \(S_z\):
\begin{equation}
[S_z,M]=0.
\end{equation}
Por tanto,
\begin{equation}
S_z\,U_D(\vk)\,S_z^{-1}=U_D\!\bigl(R_z(\pi)\vk\bigr).
\label{c25:eq:C2zD}
\end{equation}

Luego \(R_z(\pi)\) es una simetria exacta del paso Dirac minimo.

\subsection{Las otras dos rotaciones de orden dos}
El mismo razonamiento funciona para
\begin{equation}
R_x(\pi),\qquad R_y(\pi),
\end{equation}
con
\begin{equation}
S_x=\bbI_{\tau}\otimes(\ii\sigma_x),
\qquad
S_y=\bbI_{\tau}\otimes(\ii\sigma_y).
\end{equation}

Estas tres rotaciones satisfacen
\begin{equation}
R_x(\pi)^2=R_y(\pi)^2=R_z(\pi)^2=\bbI,
\qquad
R_x(\pi)R_y(\pi)=R_z(\pi),
\end{equation}
y generan un grupo discreto de cuatro elementos:
\begin{equation}
D_2=\{\bbI,\,R_x(\pi),\,R_y(\pi),\,R_z(\pi)\}.
\label{c25:eq:D2}
\end{equation}

\begin{figure}[h!]
\centering
\begin{tikzpicture}[>=Latex,node distance=1.8cm]
  \node[draw,rounded corners,fill=blue!8,minimum width=2.2cm,minimum height=0.9cm] (e) {\(\bbI\)};
  \node[draw,rounded corners,fill=green!8,minimum width=2.2cm,minimum height=0.9cm,right=of e] (x) {\(R_x(\pi)\)};
  \node[draw,rounded corners,fill=green!8,minimum width=2.2cm,minimum height=0.9cm,below=of e] (y) {\(R_y(\pi)\)};
  \node[draw,rounded corners,fill=green!8,minimum width=2.2cm,minimum height=0.9cm,below=of x] (z) {\(R_z(\pi)\)};
  \draw[->,thick] (e) -- (x);
  \draw[->,thick] (e) -- (y);
  \draw[->,thick] (x) -- node[right] {\small \(\times R_y\)} (z);
  \draw[->,thick] (y) -- node[below] {\small \(\times R_x\)} (z);
\end{tikzpicture}
\caption{En el paso minimo ordenado aparece de manera limpia un subgrupo puntual residual \(D_2\).}
\end{figure}

\subsection{Por que la rotacion de cuarto de vuelta no es exacta aqui}
Una rotacion \(R_z(\pi/2)\) envia
\begin{equation}
(k_x,k_y,k_z)\mapsto(k_y,-k_x,k_z).
\end{equation}
En el Hamiltoniano efectivo eso no trae problemas, pero en el paso exacto ocurre algo mas delicado:
\[
U_zU_yU_x
\quad\longrightarrow\quad
U_z\,U_{-x}\,U_y.
\]

Eso ya no es el mismo producto ordenado. Para restaurarlo haria falta una reorganizacion microscopica adicional del paso. Por eso, en el modelo minimo, la rotacion de \(\pi/2\) es \textbf{solo efectiva}, no exacta.

\begin{figure}[h!]
\centering
\begin{tikzpicture}[>=Latex,node distance=1.5cm]
  \node[draw,rounded corners,fill=blue!8,minimum width=4cm,minimum height=0.95cm] (a) {\(U_z\,U_y\,U_x\)};
  \node[draw,rounded corners,fill=red!8,minimum width=4cm,minimum height=0.95cm,right=4.5cm of a] (b) {\(U_z\,U_{-x}\,U_y\)};
  \draw[->,thick] (a) -- node[above,align=center] {\small rotacion de \(\pi/2\)\\\small cambia ejes y orden} (b);
\end{tikzpicture}
\caption{La memoria del orden de subpasos rompe la simetria rotacional exacta amplia del modelo minimo.}
\end{figure}

\section{Grupo microscopico del paso minimo}
Con lo anterior, la propuesta mas honesta para el grupo exacto del modelo minimo es
\begin{equation}
G_{\mathrm{micro}}=
\mathbb{Z}_t\times\bigl(\mathbb{Z}^3\rtimes D_2\bigr),
\label{c25:eq:Gmicro}
\end{equation}
donde:
\begin{enumerate}[label=\arabic*.]
    \item \(\mathbb{Z}_t\) describe los pasos enteros de tiempo;
    \item \(\mathbb{Z}^3\) describe las traslaciones espaciales de red;
    \item \(D_2\) es el subgrupo puntual residual que preserva exactamente el paso ordenado.
\end{enumerate}

Esta formula no debe sobreinterpretarse: corresponde al \textbf{modelo minimo} que venimos construyendo. Si uno simetriza el protocolo temporal o amplía la celda unitaria, el grupo puntual exacto podria crecer. Pero en esta version, el mensaje importante es otro:
\begin{displayquote}
\textbf{el grupo exacto del protoespacio discreto no es Poincare, ni tampoco el grupo cubico completo.}
\end{displayquote}

\section{Tercer nivel: el grupo efectivo infrarrojo}
\subsection{Recuperacion de rotaciones continuas}
Cuando \(a|\vp|\ll 1\), el paso se aproxima por el Hamiltoniano Dirac
\[
H_D=
v\,\tau_z\otimes(\vsigma\cdot\vp)
+m\,\tau_x\otimes\bbI.
\]

En ese nivel, el generador total de rotaciones es
\begin{equation}
\bm{J}=\bm{L}+\bm{S},
\qquad
\bm{S}=\frac{1}{2}\,\bbI_{\tau}\otimes\bm{\sigma},
\label{c25:eq:J}
\end{equation}
y se verifica
\begin{equation}
[J_i,H_D]=0.
\label{c25:eq:JH}
\end{equation}

La razon es la misma que en la ecuacion de Dirac usual: el momento \(\vp\) y el espin interno \(\vsigma\) rotan juntos y dejan invariante el escalar \(\vsigma\cdot\vp\).

Luego el grupo de rotaciones efectivas en el infrarrojo ya no es \(D_2\), sino
\begin{equation}
SO(3),
\end{equation}
con su accion espinorial levantada a
\begin{equation}
SU(2).
\end{equation}

\subsection{Aparicion de Lorentz y Poincare}
Si reescribimos el sector infrarrojo en forma covariante,
\begin{equation}
(\ii\gamma^\mu\partial_\mu-m)\Psi=0,
\end{equation}
entonces la simetria efectiva relevante deja de ser solo rotacional y pasa a ser
\begin{equation}
\mathcal{P}=SO^+(3,1)\ltimes\mathbb{R}^{3,1},
\end{equation}
es decir, el grupo de Poincare.

Al nivel espinorial, esto se levanta al recubrimiento doble
\begin{equation}
SL(2,\mathbb{C})\longrightarrow SO^+(3,1).
\end{equation}

Pero ahora la distincion metodologica es crucial:
\begin{enumerate}[label=\arabic*.]
    \item \(SO^+(3,1)\) y Poincare no son simetrias exactas del paso discreto global;
    \item son simetrias del \textbf{sector efectivo linealizado};
    \item aparecen porque el infrarrojo ya no resuelve la granularidad del protocolo microscopico.
\end{enumerate}

\section{Por que los boosts no pueden ser exactos microscopicamente}
Este punto conviene decirlo sin rodeos. Un boost continuo mezcla espacio y tiempo. Pero en el nivel exacto del modelo tenemos:
\begin{enumerate}[label=\arabic*.]
    \item espacio discreto \(\mathbb{Z}^3\);
    \item tiempo discreto \(\mathbb{Z}\);
    \item un paso elemental con orden fijo;
    \item zona de Brillouin compacta y cuasienergia periodica.
\end{enumerate}

Ese conjunto de rasgos no es compatible con un grupo continuo exacto de boosts de Lorentz. Por eso la estructura relativista solo puede emerger localmente, cerca de un nodo y para momentos pequenos.

\begin{figure}[h!]
\centering
\begin{tikzpicture}[>=Latex,node distance=1.2cm]
  \node[draw,rounded corners,fill=blue!8,minimum width=4.4cm,minimum height=0.95cm,align=center] (a) {nivel exacto\\\(\mathbb{Z}_t\times(\mathbb{Z}^3\rtimes D_2)\)};
  \node[draw,rounded corners,fill=green!8,minimum width=4.4cm,minimum height=0.95cm,align=center,below=of a] (b) {nivel infrarrojo\\\(\mathbb{R}^3\rtimes SO(3)\)};
  \node[draw,rounded corners,fill=yellow!12,minimum width=4.4cm,minimum height=0.95cm,align=center,below=of b] (c) {nivel covariante efectivo\\Poincare y \(SL(2,\mathbb{C})\)};
  \draw[->,thick] (a) -- node[right] {\small coarse graining} (b);
  \draw[->,thick] (b) -- node[right] {\small linealizacion covariante} (c);
\end{tikzpicture}
\caption{La jerarquia correcta de grupos en el proyecto.}
\end{figure}

\section{Lectura para el protoespacio discreto}
Con todo lo anterior, el objeto discreto candidato queda mejor descrito como
\begin{equation}
\mathfrak{P}_{\mathrm{disc}}=
\bigl(\Lambda,\mathbb{C}^4,U_D,G_{\mathrm{micro}}\bigr),
\label{c25:eq:proto}
\end{equation}
donde:
\begin{enumerate}[label=\arabic*.]
    \item \(\Lambda=\mathbb{Z}^3\) es la red espacial;
    \item \(\mathbb{C}^4\) es la fibra local tipo Dirac;
    \item \(U_D\) es la ley dinamica elemental;
    \item \(G_{\mathrm{micro}}\) es el grupo exacto discreto que realmente preserva la construccion.
\end{enumerate}

Esta descripcion es mejor que decir solo \enquote{la red cubica es el protoespacio}. La red sola no basta. Lo que importa es la \textbf{red mas su fibra interna mas su dinamica local mas sus simetrias exactas}.

\section{Tabla de jerarquia}
\begin{center}
\begin{tabular}{@{}lll@{}}
\toprule
Nivel & Grupo relevante & Significado \\
\midrule
Microscopico exacto & \(\mathbb{Z}_t\times(\mathbb{Z}^3\rtimes D_2)\) & paso discreto real \\
Infrarrojo espacial & \(\mathbb{R}^3\rtimes SO(3)\) & isotropia efectiva \\
Espinorial infrarrojo & \(SU(2)\) & rotaciones de espin \\
Covariante efectivo & \(SO^+(3,1)\ltimes\mathbb{R}^{3,1}\) & Dirac relativista \\
Recubrimiento espinorial & \(SL(2,\mathbb{C})\) & accion sobre espinores \\
\bottomrule
\end{tabular}
\end{center}

\section{Conclusion}
Este paso aclara una confusion muy facil de cometer. El hecho de que el sector infrarrojo reproduzca Dirac en \(3+1\) \textbf{no} significa que el modelo microscopico tenga ya el grupo de Poincare como simetria exacta.

Lo que si significa es algo mas preciso y mas valioso:
\begin{enumerate}[label=\arabic*.]
    \item existe una microdinamica discreta bien definida;
    \item esa dinamica tiene un grupo exacto discreto identificable;
    \item en el infrarrojo emerge una simetria continua mas grande;
    \item y en ese limite aparece la estructura espinorial relativista correcta.
\end{enumerate}

Eso ordena mejor la investigacion. La pregunta ya no es \enquote{si hay Dirac}, porque eso ya quedo construido. La pregunta pasa a ser:
\begin{displayquote}
\textbf{que condiciones adicionales tendria que satisfacer el grupo microscopico exacto para que la emergencia relativista no sea solo local, sino geometricamente mas robusta?}
\end{displayquote}

Ese es el siguiente paso natural del proyecto.
